\subsection{Per transmitted position, the latent message is worth less than text}
\label{sec:bandwidth}
Hidden states are often assumed to carry more than the tokens they sit on. For this receiver and these bridges, that is not what we measure. At full length, the receiver-state reference delivers 79\% of the content gain that text delivers (1.06 vs 1.33 nats; Table~\ref{tab:comm}); one difference between the two is that with the message, $B$'s first layer processes the suffix without the prefix. When only the last $K$ prefix positions are sent (Figure~\ref{fig:comm}c), $K$ tokens of text are worth more than $K$ bridged positions at every tested $K\in\{1,2,4,\dots,128\}$: the difference between text and the real-text bridge ranges from 0.08 nats [0.04, 0.12] at $K=1$ to 0.36 [0.30, 0.43] at $K=64$, and all eight Bonferroni-adjusted intervals exclude zero (smallest lower bound 0.02; intervals resample articles and are conditional on the saved bridges and one other-article donor draw). This is an operational result for this receiver, not a statement about information density. Short messages also change other things: the bridges were fitted only on 128-position messages, a truncated message places a mid-document state at sequence position 0, and absolute cross-entropy with a short latent message is \emph{worse} than with no context (Appendix~\ref{app:bandwidth}). Bridges trained for short messages, or receivers that inject at several layers or into key--value caches \citep{fu2026c2c}, may behave differently.

Equal numbers of positions are not equal communication costs. Payload: a 128-position message costs 262\,kB in float32 if $A$'s 512-dimensional states are sent and the bridge is applied at the receiver (131\,kB in float16), or 393\,kB if the bridged 768-dimensional states are sent; 128 token ids packed as 16-bit integers cost 256 bytes, so the latent message is about three orders of magnitude larger. Computation: in an idealised split implementation the receiver still runs 11 of its 12 layers on every message position, roughly the same as reading the prefix as text (a dense-layer estimate of $2\times85$M non-embedding parameters $\times128$ positions $\approx 22$ GFLOP either way), and the sender adds two of its layers; our implementation runs both models in full and overwrites states, so it does not realise even that. We did not measure latency. This design therefore shows no communication-efficiency advantage; the case for latent channels rests on settings where the sender would otherwise have to \emph{generate} text token by token, the motivation of \citet{fu2026c2c} and \citet{zou2025latentmas}, which we do not test.
